% =====================================================================
%  ГУРАВ. МАРКЕТИНГИЙН ТӨЛӨВЛӨГӨӨ
% =====================================================================
\chapter{МАРКЕТИНГИЙН ТӨЛӨВЛӨГӨӨ}

\section{Маркетингийн зорилго, зорилтот зах зээл}

\textbf{Маркетингийн зорилго:} Ажиллагаатай кофе шопуудад NFC стикер
болон ухаалаг захиалгын системээс бүрдсэн бүтээгдэхүүнийг шууд
борлуулж, нэгж борлуулалт бүрээс ашиг олох.

\textbf{Маркетингийн зорилтууд:}

\begin{itemize}
  \item Нэг жилийн дотор 10 -- 20 кофе шопт бүтээгдэхүүнийг
        борлуулж, суурилуулах.
  \item Кофе шоп тус бүрд ширээ бүрт нэг удаа NFC стикер байршуулж,
        цэс болон системийг тохируулан өгснөөр ханган нийлүүлэлтийн
        ажлыг бүрэн дуусгах.
  \item Бүтээгдэхүүнийг ашиглаж эхэлсэн кофе шопуудын санал
        хүсэлтэд үндэслэн бүтээгдэхүүнийг байнгын сайжруулах.
\end{itemize}

\textbf{Зорилтот зах зээл:} Эхний үе шатанд Улаанбаатар хотын
өндөр урсгалтай (high traffic), залуу үеийн хэрэглэгчдэд чиглэсэн
кофе шопууд. Дараагийн үе шатанд хөдөө орон нутгийн томоохон
хотууд, улмаар экспортын зах зээл.

\section{Маркетингийн хольц буюу 4P-ийн шинжилгээ}

\subsection{Product буюу Бүтээгдэхүүн}

Борлуулж буй бүтээгдэхүүн нь дараах бүрэлдэхүүн хэсгүүдээс
бүрдэнэ:

\begin{itemize}
  \item \textbf{NFC стикер сет:} Кофе шопын ширээ бүрд нэг стикер --
        тухайн ширээний дугаарыг кодчилсон, элэгдэлд тэсвэртэй
        NTAG213/215 чиптэй стикер.
  \item \textbf{Ухаалаг захиалгын систем:} Хэрэглэгчийн PWA цэс,
        захиалгын урсгал, цахим төлбөрийн интерфэйс.
  \item \textbf{Тохиргоо, нэвтрүүлэг:} Кофе шопын цэсийг системд
        бүрэн тохируулан оруулж, ажилтнуудад нь зааж өгөх.
  \item \textbf{Админ эрх:} Кофе шоп дараа нь өөрийн админ эрхээр
        нэвтрэн цэсээ шинэчлэх -- шинэ бүтээгдэхүүн нэмэх, хасах,
        үнийн өөрчлөлт оруулах боломжтой.
\end{itemize}

\textbf{Бүтээгдэхүүний нэр хүнд, чанарын баталгаа:} Системийн
таталтын ажиллагааг 99.5\%-ийн түвшинд хангах, стикерийн наалдац
болон чипийн ажиллагаанд 1 жилийн баталгаа өгөх.

\textbf{Чанарын удирдлага, хяналтын систем:} Нэгж борлуулалт
бүрийн дараа кофе шопоос хүлээн авсан санал хүсэлтийг бүртгэж,
бүтээгдэхүүний дараагийн хувилбарыг бэлтгэхэд ашиглана.

\subsection{Price буюу Үнийн бодлого}

B2B хэрэглэгчид (кофе шопууд) зориулсан нэг удаагийн багц үнэ:

\begin{table}[h]
  \centering
  \caption{Багцын үнийн бүтэц (нэг удаагийн төлбөр)}
  \label{tab:ch3-unelgee}
  \small
  \begin{tabular}{@{}p{3.2cm}p{3.0cm}p{6.6cm}@{}}
    \toprule
    \textbf{Багц} & \textbf{Үнэ} & \textbf{Багцад орсон бүтээгдэхүүн} \\
    \midrule
    Үндсэн багц & 3{,}500{,}000 \textcurrency{} & NFC стикер сет (ширеэ
    бүрд 1), захиалгын систем, цэсийг тохируулан оруулах, админ эрх,
    ажилтны сургалт \\
    Нэмэлт стикер & 50{,}000 \textcurrency{}/ширээ & Сүлжээ өргөтгөх,
    ширээ нэмэх үед стикер + тохиргоо \\
    Жилийн дэмжлэг (сонголтоор) & 600{,}000 \textcurrency{}/жил &
    Техникийн дэмжлэг, системийн шинэчлэлт, нөөц стикер 5 ширээ \\
    \bottomrule
  \end{tabular}
\end{table}

\noindent Тусгай үнэ: анхны 5 кофе шопт 20\%-ийн хөнгөлөлт
(2{,}800{,}000\textcurrency{}) өгөх «Анхдагч хэрэглэгч» урамшуулал.
Төлбөрийн комисс (QPay, SocialPay, карт) нь төлбөрийн үйлчилгээ
үзүүлэгчийн тарифаар шууд төлөгдөнө.

\subsection{Place буюу Хуваарилалтын суваг}

\begin{itemize}
  \item \textbf{Шууд борлуулалт:} Кофе шопын эзэмшигчидтэй шууд
        уулзаж, бүтээгдэхүүнийг танилцуулж, демо үзүүлэх.
  \item \textbf{Түншлэл:} Кофеийн дүн зөвлөх, тоног төхөөрөмж
        нийлүүлэгч, кафе-ийн дизайнеруудтай хамтран ажиллах.
  \item \textbf{Цахим суваг:} Веб сайт, дэлгэрэнгүй танилцуулга,
        демо туршилтын ажиллагааны систем.
\end{itemize}

\subsection{Promotion буюу Урамшуулал, идэвхжүүлэлт}

\begin{itemize}
  \item Анхдагч хэрэглэгч кофе шопуудтай хамтран түүвэр
        сурталчилгаа хийх.
  \item Кофе шопын нүүр цэс, кассын дэргэд «NFC стикертэй ширээ»
        тэмдэглэгээг ашиглах.
  \item Хэрэглэгчдэд анхны NFC уншилтаараа урамшуулал (үнэгүй
        ундаа, ундааны хөнгөлөлт) өгөхийг кофе шопуудад санал болгох.
  \item Хэрэглэгчийн эхний захиалгад 500 -- 1{,}000\textcurrency{}
        хүртэлх промо код.
\end{itemize}

\section{Борлуулалтын хүлээгдэж буй үр дүн}

Эхний жилд 10 -- 20 кофе шопт бүтээгдэхүүнийг борлуулж, дунджаар
250 ширээнд NFC стикерийг байршуулан, системийг тохируулан өгөх
зорилттой. Нэг удаагийн багцын орлого нь төслийн үндсэн орлогын
эх үүсвэр бөгөөд нэгж борлуулалт бүр шинэ кофе шоп руу гүйлгэх
сүлжээг өргөтгөнө.
